\section{Berdyaev and the Mirage of Progress}

In these passages, taken from \emph{The Destiny of Man}, \textbf{Nikolai Berdyaev} reveals his origins from Russian Hermetism and that he is more than a ``philosopher" in the academic sense. He rejects the notion that we can study man as he is today and then extrapolate back to his origins. Like \textbf{Rene Guenon}, he sees what he calls modern savages as a degenerate form and not as the proto-man. Science shows us nothing about origins, since it merely masks ``philosophic assumptions" as scientific fact. Interestingly enough, he proposes that the world was more ``plastic" than it is now. Of course, this is identical to the claims of Tradition. By ``Akashic Records", we recognize Traditional Method\footnote{\url{https://gornahoor.net/?p=624}} or the technique of Hermetic meditation\footnote{See Section \ref{sec:HermeticMeditation} in this book.}.

\begin{quotex}
Anthropologists and sociologists have devoted a great deal of attention to the primitive man but their methods and principles of investigation were determined by the evolutionary theory of the second half of the nineteenth century. They studied modern savages and from them drew conclusions about the primitive man. Scientific investigation in the strict sense was from the nature of the case impossible, but as a result of philosophic assumptions it was believe that, to being with, man was at a savage, half-animal stage and then, up to the nineteenth century, he gradually progressed. Man's distant past was inferred from his present, from savages and animals. The scientists' imagination was so poor that in man's distant past they could conceive of nothing different from what they found in modern times at the lower stages of life. But ancient man and his life were infinitely more significant and mysterious than anthropologists and sociologists suppose. In this respect theosophists and occultists are nearer the truth. There is something to be said for the Akashic Records, the Chronicle of the world, though the idea is easily vulgarized. At the dawn of humanity the world was at a different stage than it is now. It was more plastic, and the limits which divide this world from other worlds were less sharply marked. We are told this, in a covert form, in the book of Genesis. 

\end{quotex}
Berdyaev continues with his critique of scientific evolutionism and denies that the study of primitive people of today or of animal behavior can tell us anything about the mentality of ancient civilizations. Like Guenon, he denies a unitary notion of progress. On the contrary, civilizations arise and die. The myth of Atlantis may serve as a warning to us. Rather than being the pinnacle of human evolution, modern man is more likely the degenerate, fallen, and weaker product of something much greater.

\begin{quotex}
The evolutionary theory of the nineteenth century has been disproved both by science and philosophy, and cannot be used as the basis of the methods and principles of inquiry. It is inadmissible to transfer to the ancient, primitive humanity our habits of thought and feeling and our view of the world. Everything then was different, not at all similar either to the savages or to the animal world of our day. Levy Bruhl, criticizing Taylor and Fraser, tries to discover the nature of primitive thought, quite different from civilized people's thought, but his modern positivist and rationalistic mentality prevents him from understanding it. What he calls \emph{la loi de participation} shows that primitive thought was of a higher type than that of the nineteenth-century man for it expressed the mystical nearness of the knower to his object. Man loses as well as gains through the growth of civilization. He not only progresses but degenerates, falls, grows weaker and poorer. There is no doubt that some ancient knowledge connected with the proximity to the sources of being was lost by man in the course of time, and only a memory of it is left to him. There is no doubt that there existed great civilizations in the past, such as those of Babylon or Egypt, and that their fall meant a period of regress and a loss of tremendous achievements. There are considerable reasons to believe in the truth of the myth about Atlantis, where a very high civilization became morally degenerate and perished. It is far more likely that the savages as we know them are a product of degeneration and retrogression, and do not represent the primary stage of human development. In speaking of the primitive moral consciousness as we know it, we must not draw conclusions with regard to the first origins of mankind. The facts that lend themselves to study and observation are chronologically secondary and not primary. 

\end{quotex}


\flrightit{Posted on 2010-08-22 by Cologero }
